\chapter{Специфікація та сертифікація}

\section{Законодавча база}

\subsection{Загальні положення}

Базова версія <<ERP/1: Документи>> керується наступними загальними положеннями які
виражені законами України:
\\
\\
--- {2657-XII}, Про інформацію\footnote{\url{https://zakon.rada.gov.ua/laws/show/2657-XII}}, \\
--- {74\/98-ВР}, Про Національну програму інформатизації\footnote{\url{https://zakon.rada.gov.ua/laws/show/74/98-вр}}, \\
--- {393\/96-ВР}, Про звернення громадян\footnote{\url{https://zakon.rada.gov.ua/laws/show/393/96-вр}}, \\
--- {2939-VI}, Про доступ до публічної інформації\footnote{\url{https://zakon.rada.gov.ua/laws/show/2939-17}} \\
--- {2155-VIII}, Про електронні довірчі послуги\footnote{\url{https://zakon.rada.gov.ua/laws/show/2155-19}}, \\
--- {851-IV}, Про електронні документи та електронний документообіг\footnote{\url{https://zakon.rada.gov.ua/laws/show/851-15}}, \\

та розпорядженнями і постановами Кабінету Міністів України:
\\
\\
--- {386-2013-р}, Розпорядження КМУ \#3860-Р \footnote{\url{https://zakon.rada.gov.ua/laws/show/386-2013-р}}, \\
--- {373-2006-п}, Постанова КМУ \#373 \footnote{\url{https://zakon.rada.gov.ua/laws/show/373-2006-п}}. \\

\newpage
\subsection{Базова версія <<ERP/1: Документи>>}

Продукт <<ERP/1: Документи>> в основному базується на Поставі \#55 Кабінету Міністрів України та інших постановах КМУ:
\\
\\
--- {55-2018-п}, КМУ. Постанова \#55 Деякі питання документування управлінської діяльності \footnote{\url{https://zakon.rada.gov.ua/laws/show/55-2018-п}},\\
--- {749-2018-п}, КМУ. Постанова \#749 Про затвердження Порядку використання електронних довірчих послуг в органах державної влади, органах місцевого самоврядування, підприємствах, установах та організаціях державної форми власності\footnote{\url{https://zakon.rada.gov.ua/laws/show/749-2018-п}},\\
--- {v0144774-20}, ДП <<Український науково-дослідний і навчальний центр проблем стандартизації, сертифікації та якості". Наказ \#144 Про прийняття та скасування національних стандартів ДСТУ 4163:2020 та ДСТУ 9031:2020\footnote{\url{https://zakon.rada.gov.ua/rada/show/v0144774-20}},\\

але платформа продукту базується на наказаї Міністерства юстиції України, Міністерства цифрової трансформації,
Міністерства освіти і науки, та Законами України:
\\
\\
--- {z1854-12}, Міністерство юстиції України. Наказ \#1600\/5 Про затвердження Порядку роботи з електронними документами через систему електронної взаємодії органів виконавчої влади з використанням електронного цифрового підпису\footnote{\url{https://zakon.rada.gov.ua/laws/show/z1854-12}},\\
--- {z1039-20}, Міністерство цифрової трансформації України. Адміністрація державної служюи спеціального зв'язку та захисту інформації України. Наказ \#140\/614\footnote{\url{https://zakon.rada.gov.ua/laws/show/z1039-20}},\\
--- {z1306-11}, Міністерство освіти і науки, молоді та спорту України. Наказ \#1207 Про вимоги до форматів даних електронного документообігу в органах державної влади. Формат електронного повідомлення\footnote{\url{https://zakon.rada.gov.ua/laws/show/z1306-11}},\\
--- {z1421-14}, Міністерство юстиції України. Наказ \#1886/5 Про затвердження Порядку роботи з електронними документами у діловодстві та їх підготовки до передавання на архівне зберігання\footnote{\url{https://zakon.rada.gov.ua/laws/show/z1421-14}},
--- {851-IV}, Про електронні документи та електронний документообіг\footnote{\url{https://zakon.rada.gov.ua/laws/show/851-15}},\\
--- {80\/94-ВР}, Про захист інформації в інформаційно-телекомунікаційних системах\footnote{\url{https://zakon.rada.gov.ua/laws/show/80/94-вр}},

\newpage
\subsection{Розширення та додаткові модулі}

{\bf Розширення <<ERP/1: Судопис і Суддівство>>}
\\
\\
--- {4651-VI}, Кримінальний процесуальний кодекс України\footnote{\url{https://zakon.rada.gov.ua/laws/show/4651-17}},\\
--- {v0298905-20}, Офіс генерального прокурора. Наказ \#298 Про затвердження Положення про Єдиний реєстр досудових розслідувань, порядок його формування та ведення\footnote{\url{https://zakon.rada.gov.ua/laws/show/v0298905-20}},
\\
\\
{\bf Розширення <<ERP/1: Закупівлі>>}
\\
\\
--- {922-19}, Закон України про публічні закупівлі\footnote{\url{https://zakon.rada.gov.ua/laws/show/922-19}},
--- {169}, Постанова КМУ \#169, \\
--- {1178} Постанова КМУ \#1178, \\
--- {808-20} Закон України про оборонні закупівлі, \\
--- {1275-2022-п} Постанова КМУ \#1275, \\
--- {1070-2019-п} Постанова КМУ \#1070, \\
--- {224-2020-п} Постанова КМУ \#224, \\
--- {710-2016-п} Постанова КМУ \#710, \\
--- {z0500-20} Наказ Мінекономіки \#708, \\
--- {544-2016-п} Постанова КМУ \#544, \\
--- {v1749731-15} Наказ Мінекономіки \#1749, \\
--- {1495-п} Постанова КМУ \#1495. \\
\\
\\
{\bf Розширення <<ERP/1: Зброя>>}
\\
\\
--- {5708}, Проект Закону про право на цивільну вогнепальну зброю\footnote{\url{http://w1.c1.rada.gov.ua/pls/zweb2/webproc4_2?pf3516=5708&skl=10}},
--- {5709}, Проект Закону про внесення змін до Кодексу України про адміністративні правопорушення та Кримінального кодексу України для реалізації положень Закону України "Про право на цивільну вогнепальну зброю"\footnote{\url{http://w1.c1.rada.gov.ua/pls/zweb2/webproc4_2?pf3516=5709&skl=10}},
\\
\newpage
{\bf Розширення <<ERP/1: Військова частина>>}
\\
\\
--- {ДСТУ 2732-94}, Діловодство і архівна справа. \\
--- {26.05.2014 \#333}, Інструкція з ведення обліку особового складу. \\
--- {\#608} Порядок проведення службового розслідування. \\
--- {26.07.2018 \#370}, Інструкція з діловодства. \\
--- {07.04.2017 \#124}, Інструкція з діловодства. \\
--- {07.10.2015 \#393}, Положення Про Юридичну Службу. \\
--- {29.11.2018 \#604}, Інструкція з надання доповідей і донесень про події,
                кримінальні правопорушення, військові адміністративні
                правопорушення та адміністративні правопорушення,
                пов'язані з корупцією, порушення військової дисципліни та їх облік. \\

\section{Функціональні модулі та вимоги}

\subsection{Призначення та цілі впровадження}

\subsection{Ревізія поточних систем}

\newpage
\chapter{Політики і вимоги}

\section{Політики}

\subsection{Загальні принципи}

\section{Технічні вимоги}

\subsection{Перелік технічних вимог наведених в політиках}

\subsection{Повний зміст технічних вимог}

\section{Технічне завдання}

\subsection*{Зміст технічного завдання}

\section{Посилання на стандарти ДСТУ}

\newpage
\section{Функціональні вимоги}

\subsection{Вимоги до процесів системи}

\subsection{Вимоги до викликів API}

\subsection{Вимоги до інтерфейсу користувача}

\subsubsection{Обробка помилок}

\subsubsection{Загальні вимоги}

\newpage
\section{Нефункціональні вимоги}

\subsection{Вимоги до архітектури}

\subsubsection{Визначення термінів і призначення документа}

\subsubsection{Рівні архітектури компонентів системи}

\subsubsection{Маніфест відповідності архітектурі}

\subsubsection{Таксономія системних компонент}

\subsection{Вимоги до безпеки}

\subsubsection{Загальні вимоги до інформаційної безпеки}

\subsubsection{Технічні вимоги до інформаційної безпеки}

\subsection{Вимоги до потужності та ємності}

\subsubsection{Таблиця характеристик продуктивності}

\subsection{Вимоги до системи логування}

\subsubsection{Загальні вимоги}

\subsubsection{Рівні логування}

\subsubsection{Інформація для вендорів про інструментарії логування}

\subsubsection{Події логування}

\subsubsection{Структура журналу логування}

\subsubsection{Вимоги до структурних логів ELK стеку}

\newpage
\subsection{Вимоги до контролю якості}

\subsubsection{Загальні вимоги до обліку і контролю якості}

\subsubsection{Вимоги до ручного тестування}

\subsubsection{Вимоги до формату сценаріїв та протоколів тестування}

\subsubsection{Вимоги до плану тестування}

\subsubsection{Вимоги до автоматичного тестування}

\subsubsection{Вимоги до результатів приймального (UAT) тестування}

\subsection{Вимоги до інтерфейсів}

\subsubsection{Загальні вимоги}

\subsubsection{Вимоги до відображення інформації}

\subsubsection{Вимоги до швидкодії додатка}

\subsubsection{Вимоги до технології реалізації додатку}

\subsection{Юридичні вимоги}

\addtocontents{toc}{\protect\newpage}
\newpage
\chapter{Служби ERP/1}
\section{Шаблони документів}

\subsection{Вимоги до сутностей}

\subsubsection{Шаблон DOCX/XLSX}

\subsubsection{Набір змінних}

\subsection{Вимоги до процесів}

\subsubsection{Заповнення шаблону}

\subsubsection{Версіонування шаблонів}

\newpage
\section{Генерація документів}

\subsection{Вимоги до сутностей}

\subsubsection{Заповнений документ}

\subsubsection{PDF-документ}

\subsection{Вимоги до процесів}

\subsubsection{Генерація PDF}

\subsubsection{Накладання КЕП}

\subsubsection{Архівація}

\section{Генерація бар- та штрих-кодів}

\subsection{Вимоги до сутностей}

\subsubsection{Код}

\subsubsection{Накладання}

\subsection{Вимоги до процесів}

\subsubsection{Генерація коду}

\subsubsection{Накладання на документ}

\subsubsection{Верифікація}

\newpage
\section{Розпізнавання Tesseract}

\subsection{Вимоги до сутностей}

\subsubsection{Документ PDF}

\subsubsection{Результати OCR}

\subsection{Вимоги до процесів}

\subsubsection{Інгестія документа}

\subsubsection{Виконання OCR}

\subsubsection{Повнотекстовий пошук}

\subsubsection{Отримання підсвічування}

\newpage
\section{Система сканування}

\subsection{Вимоги до сутностей}

\subsubsection{Профіль сканера}

\subsubsection{Сканований документ}

\subsection{Вимоги до процесів}

\subsubsection{Ініціація сканування}

\subsubsection{Доставка результату}

\newpage
\section{Багатокористувацький редактор}

\subsection{Вимоги до сутностей}

\subsubsection{Редакційний документ}

\subsubsection{Операція редагування}

\subsection{Вимоги до процесів}

\subsubsection{Синхронізація в реальному часі}

\subsubsection{Історія змін}

\subsubsection{Фіксація версії}

\subsubsection{Транзакційність}

\subsection{Призначення та цілі впровадження}

\newpage
\section{Інфраструктура безпеки}

\subsection{Провайдери української криптографії}

\subsection{Вимоги до сутностей}

\subsubsection{Директорія (LDAP)}

\subsubsection{Сертифікати (X.509)}

\subsubsection{Користувачі (inetOrgPerson)}

\subsubsection{Організації (organization)}

\subsubsection{Криптографічне повідомлення (CMS X.894)}

\subsubsection{Кваліфікований електронний підпис (ДСТУ 4145)}

\subsubsection{Відкликання і перевірка (OCSP)}

\subsubsection{Мітка часу (TSP)}

\subsubsection{Ключі шифрування та сервісні креденціали}

\subsection{Вимоги до процесів}

\subsubsection{Видача сертифікату (enroll)}

\subsubsection{Автентифікація сертифікатом (check)}

\subsubsection{Відкликання сертифікату (revoke)}

\subsubsection{Оновлення сертифікату (renew)}

\subsubsection{Підпис і верифікація (sign/verify)}

\subsubsection{Шифрування і розшифрування (encrypt/decrypt)}

\subsubsection{Управління секретами (store/retrieve)}

\subsection{Вимоги до серіалізації і сховища}

\newpage
\section{Контроль доступу}

\subsection{Вимоги до рольової моделі ABAC}

\subsubsection{Об'єкт доступу}

\subsubsection{Суб'єкт доступу}

\subsubsection{Правило доступу}

\subsection{Вимоги до сутностей ABAC}

\subsubsection{Політики доступу}

\subsection{Вимоги до процесів ABAC}

\subsubsection{Життєвий цикл Політики доступу}

\subsubsection{Точки доступу}

\newpage
\section{Директорія підприємства}

\subsection{Вимоги до сутностей директорії}

\subsubsection{Адміністратор організацій}

\subsubsection{Адміністратор структурних підрозділів}

\subsubsection{Адміністратор співробітників}

\subsubsection{Адміністратор користувачів}

\subsection{Вимоги до процесів директорії}

\subsubsection{Життєвий цикл Адміністратора}

\subsubsection{Життєвий цикл Організації}

\subsubsection{Життєвий цикл Користувача}

\subsubsection{Життєвий цикл Сесії}

\subsubsection{Життєвий цикл Ролі}

\subsection{Вимоги до рольової моделі структурних одиниць}

\subsubsection{Точки доступу}

\subsection{Вимоги до сховища}

\subsection{Вимоги до серіалізації і валідації}

\newpage
\section{Оркестрація процесів}

\subsection{Вимоги до сутностей BPMN}

\subsubsection{Документ}

\subsubsection{Процес}

\subsubsection{Монітор}

\subsubsection{Задача користувача (АРМ)}

\subsubsection{Системна задача}

\subsubsection{Перехід}

\subsubsection{Подія}

\subsubsection{Таймаут}

\subsubsection{Асинхронне повідомлення}

\subsection{Вимоги до процесів}

\subsubsection{Створення}

\subsubsection{Зупинення}

\subsubsection{Відновлення}

\subsubsection{Архівація}

\subsubsection{Зміна контексту процесу}

\subsubsection{Перехід на стадію}

\subsection{Вимоги до сховища}

\newpage
\section{Словникова підсистема}

\subsection{Загальні відомості}

\subsection{Вимоги до сутностей HL7}

\subsubsection{Системноутворюючі довідники}

\subsubsection{Похідні довідники}

\subsubsection{Система імен}

\subsubsection{Відображення довідників}

\subsubsection{Налаштування довідникової системи}

\subsection{Вимоги до процесів}

\subsubsection{Створення попереднього словника}

\subsubsection{Редагування попереднього словника}

\subsubsection{Публікація словника}

\subsubsection{Деактивація словника}

\subsubsection{Архівування словника}

\subsection{Вимоги до сховища}

\section{Мета сервіси та конструктор}

\subsection{Реєстр сутностей системи (API)}

\subsection{Біблітека валідації (JSON Schema)}

\subsection{Конструктор сутностей та форм}

\subsection{Конструктор бізнес процесів}

\subsection{Конструктор словників}

\subsection{Конструктор правил доступу}

\subsection{Конструктор шаблонів профілів користувачів}

